% !TEX root = ../Document.tex
% Résumé du mémoire.
% Abstract in French.
%
\chapter*{RÉSUMÉ}\thispagestyle{headings}
\addcontentsline{toc}{compteur}{RÉSUMÉ}

L'intégration de politiques génératives, telles que le \textit{Flow Matching}, offre une solution efficace pour la génération de trajectoires réactives et multi-modales en robotique de manipulation à partir d'observations visuelles. Toutefois, ces modèles appris ne fournissent aucune garantie formelle de sécurité et peuvent commander des mouvements menant à des collisions dans des environnements contenant des obstacles non vus à l'entraînement. À l'inverse, les méthodes de filtrage réactif fondées sur les fonctions barrières de contrôle (\ac{CBF}) garantissent l'invariance formelle d'ensembles sûrs, mais s'appuient traditionnellement sur des approximations géométriques simplifiées du robot et de l'environnement.

Ce mémoire présente une architecture découplée à deux niveaux pour l'exécution sécurisée en temps réel de trajectoires générées par une politique de \textit{Flow Matching} gelée, sans réentraînement de celle-ci. Le système sépare entièrement la planification de la tâche et la garantie de sécurité. La perception réunit un modèle de segmentation sémantique (\ac{SAM}~3) et une carte de points persistante permettant de maintenir la représentation des obstacles lors des occlusions causées par les mouvements d'une caméra embarquée au poignet. La politique nominale génère à une cadence de $5$~Hz des consignes de vitesse articulaire. Un bouclier réactif temps réel évalue un champ de distance signée (\ac{SDF}) appris par polynômes de Bernstein pour le corps complet du robot et résout, à $50$~Hz par programmation quadratique, les contraintes \ac{CBF} associées. Cette couche est complétée à $1$~kHz par un frein réflexe qui ajuste dynamiquement l'horizon de certification entre les résolutions du solveur.

L'architecture est évaluée en simulation Gazebo sur deux tâches de manipulation (alimentation assistée et saisie-dépôt) avec un bras Franka Emika Panda et un bras uFactory xArm7. Les résultats expérimentaux démontrent que le filtre de sécurité élimine l'intégralité des $51$ collisions observées en régime nominal sur $60$ essais sans filtre, tout en préservant le succès des tâches ($30/30$ en saisie-dépôt) et en induisant une déviation moyenne de l'outil inférieure à $0{,}6$~cm. Le déploiement sur le second manipulateur valide le transfert du bouclier de sécurité par le seul réentraînement de la cinématique et des champs de distance par lien. Ce travail confirme qu'une séparation stricte entre la politique générative et le filtre réactif permet de concilier la flexibilité des modèles appris avec des garanties formelles de sécurité en robotique de manipulation.
